% Recovered verbatim from chapters_16_20_with_part/chapters_16_20_with_part.tex
% Canonical chapter 17
\chapter{Entanglement, ER Bridges, and Locality}

\section{ER$=$EPR in MK--43 Ultra}

The ER network is not merely a geometric scaffold; it directly encodes entanglement patterns in the underlying quantum state.
Pairs or regions of degrees of freedom with high mutual entanglement are connected by strong ER links, while weakly entangled regions have sparse or weak connectivity.

In this sense, the ER$=$EPR slogan is concretised:
\begin{quote}
  entangled pairs and regions are joined by ER bridges in the microscopic network, and the pattern of these bridges determines effective adjacency and geometry.
\end{quote}

\section{Mutual Information and ER Weights}

A natural ansatz is that the ER weights are monotone functions of mutual information:
\[
  W_{ij} \propto f\big(I(i:j)\big),
\]
where $I(i:j)$ is the mutual information between degrees of freedom associated with sites $i$ and $j$.
The exact form of $f$ is model-dependent, but the monotonicity ensures that highly entangled pairs are strongly connected in the ER graph.

\section{Locality as an Entanglement Pattern}

Locality is no longer fundamental; it is an emergent property of entanglement structure.
Regions of spacetime are ``nearby'' if their microscopic degrees of freedom are more strongly entangled with each other than with the rest of the universe.
Classical causal structure thus arises from the organisation of ER links, constrained by the dynamics of the phase and chromatic fields.

